\hypertarget{advanced-topics}{%
\chapter{ADVANCED TOPICS}\label{advanced-topics}}

\hypertarget{template-metaprogramming}{%
\section{TEMPLATE METAPROGRAMMING}\label{template-metaprogramming}}

\hypertarget{compile-time-computation}{%
\section{1.1 Compile-Time Computation}\label{compile-time-computation}}

Template metaprogramming enables computation at compile-time.

\hypertarget{factorial-at-compile-time}{%
\subsection{Factorial at Compile-Time}\label{factorial-at-compile-time}}

\begin{lstlisting}[language={C++}]
#include <iostream>
using namespace std;

// Base case
template<int N>
struct Factorial {
    static constexpr int value = N * Factorial<N-1>::value;
};

// Specialization (base case)
template<>
struct Factorial<0> {
    static constexpr int value = 1;
};

int arr[Factorial<5>::value];  // Array of size 120!

cout << Factorial<10>::value << "\n";  // 3628800
\end{lstlisting}

\hypertarget{fibonacci-at-compile-time}{%
\subsection{Fibonacci at Compile-Time}\label{fibonacci-at-compile-time}}

\begin{lstlisting}[language={C++}]
template<int N>
struct Fib {
    static constexpr int value = Fib<N-1>::value + Fib<N-2>::value;
};

template<>
struct Fib<0> {
    static constexpr int value = 0;
};

template<>
struct Fib<1> {
    static constexpr int value = 1;
};

// Memoization to avoid exponential time
template<int N>
struct FibMemo {
    static constexpr int value = FibMemo<N-1>::value + FibMemo<N-2>::value;
};

static constexpr int fib20 = FibMemo<20>::value;  // Computed at compile-time
\end{lstlisting}

\hypertarget{compile-time-power-check}{%
\subsection{Compile-Time Power Check}\label{compile-time-power-check}}

\begin{lstlisting}[language={C++}]
template<unsigned int N>
struct IsPowerOfTwo {
    static constexpr bool value = (N > 0) && ((N & (N - 1)) == 0);
};

static_assert(IsPowerOfTwo<16>::value);
static_assert(!IsPowerOfTwo<7>::value);
\end{lstlisting}

\begin{center}\rule{0.5\linewidth}{0.5pt}\end{center}

\hypertarget{template-specialization}{%
\section{1.2 Template Specialization}\label{template-specialization}}

\hypertarget{full-specialization}{%
\subsection{Full Specialization}\label{full-specialization}}

\begin{lstlisting}[language={C++}]
// Primary template
template<typename T, typename U>
struct Pair {
    static void info() { cout << "Generic pair\n"; }
};

// Full specialization
template<>
struct Pair<int, string> {
    static void info() { cout << "int-string pair\n"; }
};

// Full specialization for pointers
template<typename T>
struct Pair<T*, T*> {
    static void info() { cout << "Two pointers\n"; }
};

Pair<int, string>::info();     // "int-string pair"
Pair<double, string>::info();  // "Generic pair"
Pair<int*, int*>::info();      // "Two pointers"
\end{lstlisting}

\hypertarget{partial-specialization}{%
\subsection{Partial Specialization}\label{partial-specialization}}

\begin{lstlisting}[language={C++}]
// Primary
template<typename T, typename U>
struct Container {
    static void type() { cout << "Generic\n"; }
};

// Partial specialization - same types
template<typename T>
struct Container<T, T> {
    static void type() { cout << "Same type\n"; }
};

// Partial specialization - pointers
template<typename T, typename U>
struct Container<T*, U*> {
    static void type() { cout << "Two pointers\n"; }
};

// Partial specialization - array
template<typename T, int N>
struct Container<T[N], T> {
    static void type() { cout << "Array and element\n"; }
};

Container<int, int>::type();           // "Same type"
Container<int*, double*>::type();      // "Two pointers"
Container<int[5], int>::type();        // "Array and element"
\end{lstlisting}

\hypertarget{advanced-metaprogramming-patterns}{%
\section{1.3 Advanced Metaprogramming
Patterns}\label{advanced-metaprogramming-patterns}}

\hypertarget{typelists}{%
\subsection{Typelists}\label{typelists}}

A list of types at compile-time, essential for ECS and Variant
implementation.

\begin{lstlisting}[language={C++}]
template<typename... Ts>
struct TypeList {};

// Length of list
template<typename List> struct Length;

template<typename... Ts>
struct Length<TypeList<Ts...>> {
    static constexpr size_t value = sizeof...(Ts);
};

// Access type at index
template<size_t N, typename List> struct At;

template<typename Head, typename... Tail>
struct At<0, TypeList<Head, Tail...>> {
    using type = Head;
};

template<size_t N, typename Head, typename... Tail>
struct At<N, TypeList<Head, Tail...>> {
    using type = typename At<N-1, TypeList<Tail...>>::type;
};

// Usage
using MyTypes = TypeList<int, float, double>;
static_assert(Length<MyTypes>::value == 3);
static_assert(std::is_same_v<At<1, MyTypes>::type, float>);
\end{lstlisting}

\begin{center}\rule{0.5\linewidth}{0.5pt}\end{center}

\hypertarget{sfinae-type-traits}{%
\section{SFINAE \& TYPE TRAITS}\label{sfinae-type-traits}}

\hypertarget{sfinae---substitution-failure-is-not-an-error}{%
\section{2.1 SFINAE - Substitution Failure Is Not An
Error}\label{sfinae---substitution-failure-is-not-an-error}}

SFINAE enables overload resolution based on template parameter
substitution.

\hypertarget{basic-sfinae}{%
\subsection{Basic SFINAE}\label{basic-sfinae}}

\begin{lstlisting}[language={C++}]
#include <type_traits>

// Enable if T is integral
template<typename T>
enable_if_t<is_integral_v<T>>
process(T x) {
    cout << "Integer: " << x << "\n";
}

// Enable if T is floating point
template<typename T>
enable_if_t<is_floating_point_v<T>>
process(T x) {
    cout << "Float: " << x << "\n";
}

process(42);      // "Integer: 42"
process(3.14);    // "Float: 3.14"
\end{lstlisting}

\hypertarget{detector-pattern}{%
\subsection{Detector Pattern}\label{detector-pattern}}

\begin{lstlisting}[language={C++}]
// Detect if type has value_type
template<typename T, typename = void>
struct HasValueType : false_type {};

template<typename T>
struct HasValueType<T, void_t<typename T::value_type>> : true_type {};

// Usage
static_assert(HasValueType<vector<int>>::value);
static_assert(!HasValueType<int>::value);
\end{lstlisting}

\hypertarget{advanced-sfinae-with-multiple-conditions}{%
\subsection{Advanced SFINAE with Multiple
Conditions}\label{advanced-sfinae-with-multiple-conditions}}

\begin{lstlisting}[language={C++}]
template<typename T>
enable_if_t<
    is_copy_constructible_v<T> &&
    is_move_constructible_v<T> &&
    is_equality_comparable_v<T>
>
smart_copy(const T& src, T& dst) {
    dst = src;
}
\end{lstlisting}

\begin{center}\rule{0.5\linewidth}{0.5pt}\end{center}

\hypertarget{type-traits-mastery}{%
\section{2.2 Type Traits Mastery}\label{type-traits-mastery}}

\hypertarget{custom-type-traits}{%
\subsection{Custom Type Traits}\label{custom-type-traits}}

\begin{lstlisting}[language={C++}]
// Check if type has a begin() and end()
template<typename T, typename = void>
struct IsContainer : false_type {};

template<typename T>
struct IsContainer<T, void_t<
    decltype(declval<T>().begin()),
    decltype(declval<T>().end())
>> : true_type {};

// Check if callable with specific signature
template<typename F, typename... Args>
struct IsCallable : false_type {};

template<typename F, typename... Args>
struct IsCallable<F, 
    void_t<decltype(declval<F>()(declval<Args>()...))>
> : true_type {};

// Usage
static_assert(IsContainer<vector<int>>::value);
static_assert(!IsContainer<int>::value);

auto lambda = [](int x) { return x; };
static_assert(IsCallable<decltype(lambda), int>::value);
\end{lstlisting}

\begin{center}\rule{0.5\linewidth}{0.5pt}\end{center}

\hypertarget{expression-templates}{%
\section{EXPRESSION TEMPLATES}\label{expression-templates}}

\hypertarget{lazy-evaluation-pattern}{%
\section{3.1 Lazy Evaluation Pattern}\label{lazy-evaluation-pattern}}

Expression templates defer computation until needed.

\hypertarget{vector-operations}{%
\subsection{Vector Operations}\label{vector-operations}}

\begin{lstlisting}[language={C++}]
#include <vector>
#include <iostream>

// Expression template for vector operations
template<typename Expr>
class VectorExpr {
public:
    double operator[](int i) const {
        return static_cast<const Expr&>(*this)[i];
    }
    
    int size() const {
        return static_cast<const Expr&>(*this).size();
    }
};

class Vector : public VectorExpr<Vector> {
private:
    std::vector<double> data;
    
public:
    Vector(int n) : data(n) {}
    
    double operator[](int i) const { return data[i]; }
    int size() const { return data.size(); }
    double& operator[](int i) { return data[i]; }
};

// Addition expression (no computation yet)
template<typename L, typename R>
class VectorSum : public VectorExpr<VectorSum<L, R>> {
private:
    const L& lhs;
    const R& rhs;
    
public:
    VectorSum(const L& l, const R& r) : lhs(l), rhs(r) {}
    
    double operator[](int i) const {
        return lhs[i] + rhs[i];
    }
    
    int size() const { return lhs.size(); }
};

// Operator overloading creates expression tree
template<typename L, typename R>
VectorSum<L, R> operator+(const VectorExpr<L>& l, const VectorExpr<R>& r) {
    return VectorSum<L, R>(static_cast<const L&>(l), static_cast<const R&>(r));
}

// Scalar multiplication
template<typename Expr>
class VectorScale : public VectorExpr<VectorScale<Expr>> {
private:
    const Expr& expr;
    double scale;
    
public:
    VectorScale(const Expr& e, double s) : expr(e), scale(s) {}
    
    double operator[](int i) const {
        return expr[i] * scale;
    }
    
    int size() const { return expr.size(); }
};

template<typename Expr>
VectorScale<Expr> operator*(double s, const VectorExpr<Expr>& e) {
    return VectorScale<Expr>(static_cast<const Expr&>(e), s);
}

// Usage
int main() {
    Vector v1(3), v2(3), v3(3);
    v1[0] = 1; v1[1] = 2; v1[2] = 3;
    v2[0] = 4; v2[1] = 5; v2[2] = 6;
    
    // No computation yet!
    auto expr = v1 + v2 + 2.0 * v3;
    
    // Computation happens here
    for (int i = 0; i < 3; i++) {
        cout << expr[i] << " ";
    }
    
    return 0;
}
\end{lstlisting}

\hypertarget{triggering-computation-assignment}{%
\section{3.2 Triggering Computation
(Assignment)}\label{triggering-computation-assignment}}

To make \passthrough{\lstinline!v3 = v1 + v2!} work efficiently, we add
an assignment operator to \passthrough{\lstinline!Vector!} that takes
any \passthrough{\lstinline!VectorExpr!}.

\begin{lstlisting}[language={C++}]
// Inside class Vector
template <typename Expr>
Vector& operator=(const VectorExpr<Expr>& expr) {
    if (size() != expr.size()) {
        // resize...
    }
    for (int i = 0; i < size(); ++i) {
        data[i] = expr[i]; // Evaluates tree at index i
    }
    return *this;
}
\end{lstlisting}

\textbf{Why is this fast?} It expands to:
\passthrough{\lstinline!v3[i] = v1[i] + v2[i]!}. There are \textbf{zero
temporary vectors} created. No allocations. Just one loop.

\begin{center}\rule{0.5\linewidth}{0.5pt}\end{center}

\hypertarget{policy-based-design}{%
\section{POLICY-BASED DESIGN}\label{policy-based-design}}

\hypertarget{template-policies}{%
\section{4.1 Template Policies}\label{template-policies}}

Policy-based design separates concerns using template parameters.

\hypertarget{thread-safety-policy}{%
\subsection{Thread Safety Policy}\label{thread-safety-policy}}

\begin{lstlisting}[language={C++}]
#include <mutex>
#include <memory>

// Policy: No synchronization
struct NoSync {
    void lock() {}
    void unlock() {}
};

// Policy: Mutex-based
struct MutexSync {
private:
    mutable std::mutex m;
    
public:
    void lock() { m.lock(); }
    void unlock() { m.unlock(); }
};

// Generic container with synchronization policy
template<typename T, typename SyncPolicy = NoSync>
class SafeVector : private SyncPolicy {
private:
    std::vector<T> data;
    
public:
    void push_back(const T& value) {
        this->lock();
        data.push_back(value);
        this->unlock();
    }
    
    T pop_back() {
        this->lock();
        T value = data.back();
        data.pop_back();
        this->unlock();
        return value;
    }
};

// Usage
SafeVector<int> single_threaded;  // No synchronization
SafeVector<int, MutexSync> multi_threaded;  // Thread-safe
\end{lstlisting}

\hypertarget{storage-policy}{%
\subsection{Storage Policy}\label{storage-policy}}

\begin{lstlisting}[language={C++}]
// Policy: Dynamic allocation
struct DynamicStorage {
    template<typename T>
    using Allocator = std::allocator<T>;
};

// Policy: Static allocation
template<int MaxSize>
struct StaticStorage {
    // Allocate from stack
};

template<typename T, typename StoragePolicy>
class Container : private StoragePolicy {
private:
    std::vector<T> data;
};
\end{lstlisting}

\hypertarget{policy-based-smart-pointer-advanced-example}{%
\section{4.2 Policy-Based Smart Pointer (Advanced
Example)}\label{policy-based-smart-pointer-advanced-example}}

A smart pointer is defined by: 1. \textbf{Storage}: How it stores the
pointer (raw vs compressed). 2. \textbf{Ownership}: Ref-counted vs
Unique vs Linked. 3. \textbf{Checking}: Assert on access vs No check.

\begin{lstlisting}[language={C++}]
template <
    class T,
    template <class> class CheckingPolicy,
    template <class> class ThreadingPolicy
>
class SmartPtr : public CheckingPolicy<T>, public ThreadingPolicy<T> {
    T* ptr;
public:
    T* operator->() {
        CheckingPolicy<T>::check(ptr);
        ThreadingPolicy<T>::lock(); // Fake lock example
        return ptr;
    }
};
\end{lstlisting}

This approach allows generating thousands of smart pointer variants with
zero runtime overhead (all resolved at compile-time).

\begin{center}\rule{0.5\linewidth}{0.5pt}\end{center}

\hypertarget{memory-management-optimization}{%
\section{MEMORY MANAGEMENT \&
OPTIMIZATION}\label{memory-management-optimization}}

\hypertarget{custom-allocators}{%
\section{5.1 Custom Allocators}\label{custom-allocators}}

\begin{lstlisting}[language={C++}]
#include <memory>

template<typename T>
class ArenaAllocator {
private:
    static constexpr size_t ARENA_SIZE = 10000;
    char arena[ARENA_SIZE];
    char* current;
    
public:
    ArenaAllocator() : current(arena) {}
    
    T* allocate(size_t n) {
        if (current + n * sizeof(T) > arena + ARENA_SIZE) {
            throw std::bad_alloc();
        }
        T* ptr = reinterpret_cast<T*>(current);
        current += n * sizeof(T);
        return ptr;
    }
    
    void deallocate(T* p, size_t n) {
        // No deallocation for arena allocator
    }
};

// Usage
vector<int, ArenaAllocator<int>> fast_vector;
\end{lstlisting}

\hypertarget{small-object-optimization-soo}{%
\section{5.2 Small Object Optimization
(SOO)}\label{small-object-optimization-soo}}

\begin{lstlisting}[language={C++}]
template<typename T, size_t SmallSize = 32>
class SmallVector {
private:
    union {
        T* ptr;              // Dynamic allocation
        std::aligned_storage_t<SmallSize, alignof(T)> small;
    };
    
    size_t size_val;
    bool is_small;
    
public:
    SmallVector() : size_val(0), is_small(true) {}
    
    void push_back(const T& value) {
        if (size_val < SmallSize / sizeof(T)) {
            // Use small storage
            new (&small) T(value);
            is_small = true;
        } else {
            // Switch to dynamic
            if (is_small) {
                // Copy small to dynamic
                ptr = new T[size_val + 1];
                is_small = false;
            }
            ptr[size_val] = value;
        }
        size_val++;
    }
    
    ~SmallVector() {
        if (!is_small) {
            delete[] ptr;
        }
    }
};
\end{lstlisting}

\begin{center}\rule{0.5\linewidth}{0.5pt}\end{center}

\hypertarget{concurrency-parallelism}{%
\section{CONCURRENCY \& PARALLELISM}\label{concurrency-parallelism}}

\hypertarget{advanced-threading-patterns}{%
\section{6.1 Advanced Threading
Patterns}\label{advanced-threading-patterns}}

\hypertarget{thread-pool}{%
\subsection{Thread Pool}\label{thread-pool}}

\begin{lstlisting}[language={C++}]
#include <thread>
#include <queue>
#include <mutex>
#include <condition_variable>
#include <functional>

class ThreadPool {
private:
    vector<std::thread> workers;
    queue<std::function<void()>> tasks;
    mutex task_mutex;
    condition_variable cv;
    bool stop = false;
    
public:
    ThreadPool(size_t num_threads) {
        for (size_t i = 0; i < num_threads; i++) {
            workers.emplace_back([this] {
                while (true) {
                    std::function<void()> task;
                    
                    {
                        unique_lock<mutex> lock(task_mutex);
                        cv.wait(lock, [this] { return !tasks.empty() || stop; });
                        
                        if (stop && tasks.empty()) return;
                        
                        task = std::move(tasks.front());
                        tasks.pop();
                    }
                    
                    task();
                }
            });
        }
    }
    
    template<typename F, typename... Args>
    auto enqueue(F&& f, Args&&... args) {
        using return_type = std::invoke_result_t<F, Args...>;
        
        auto task = make_shared<std::packaged_task<return_type()>>(
            bind(forward<F>(f), forward<Args>(args)...)
        );
        
        auto result = task->get_future();
        
        {
            unique_lock<mutex> lock(task_mutex);
            tasks.emplace([task] { (*task)(); });
        }
        
        cv.notify_one();
        return result;
    }
    
    ~ThreadPool() {
        {
            unique_lock<mutex> lock(task_mutex);
            stop = true;
        }
        cv.notify_all();
        for (auto& worker : workers) {
            worker.join();
        }
    }
};

// Usage
ThreadPool pool(4);

auto future = pool.enqueue([](int x) { return x * 2; }, 42);
cout << future.get() << "\n";  // 84
\end{lstlisting}

\hypertarget{lock-free-queue}{%
\subsection{Lock-Free Queue}\label{lock-free-queue}}

\begin{lstlisting}[language={C++}]
#include <atomic>

template<typename T>
class LockFreeQueue {
private:
    struct Node {
        T value;
        atomic<Node*> next{nullptr};
    };
    
    atomic<Node*> head;
    atomic<Node*> tail;
    
public:
    LockFreeQueue() {
        auto dummy = new Node();
        head.store(dummy, memory_order_relaxed);
        tail.store(dummy, memory_order_relaxed);
    }
    
    void push(const T& value) {
        auto new_node = new Node{value};
        Node* old_tail = tail.load(memory_order_acquire);
        
        old_tail->next.store(new_node, memory_order_release);
        tail.store(new_node, memory_order_release);
    }
    
    bool pop(T& value) {
        Node* old_head = head.load(memory_order_acquire);
        Node* next = old_head->next.load(memory_order_acquire);
        
        if (next == nullptr) return false;
        
        value = next->value;
        head.store(next, memory_order_release);
        delete old_head;
        
        return true;
    }
};
\end{lstlisting}

\begin{center}\rule{0.5\linewidth}{0.5pt}\end{center}

\hypertarget{type-erasure-patterns}{%
\section{TYPE ERASURE PATTERNS}\label{type-erasure-patterns}}

\hypertarget{virtual-function-based-type-erasure}{%
\section{7.1 Virtual Function-Based Type
Erasure}\label{virtual-function-based-type-erasure}}

\begin{lstlisting}[language={C++}]
class AnyCallable {
private:
    struct Base {
        virtual ~Base() = default;
        virtual void call() = 0;
    };
    
    template<typename F>
    struct Derived : Base {
        F func;
        Derived(F f) : func(f) {}
        void call() override { func(); }
    };
    
    unique_ptr<Base> impl;
    
public:
    template<typename F>
    AnyCallable(F f) : impl(make_unique<Derived<F>>(f)) {}
    
    void operator()() {
        impl->call();
    }
};

// Usage
AnyCallable c1([](){ cout << "Lambda\n"; });
AnyCallable c2(std::bind(...));
AnyCallable c3(&function);

c1();  // "Lambda"
\end{lstlisting}

\hypertarget{stdfunction-implementation}{%
\section{7.2 std::function
Implementation}\label{stdfunction-implementation}}

\begin{lstlisting}[language={C++}]
#include <memory>

template<typename>
class Function;

template<typename R, typename... Args>
class Function<R(Args...)> {
private:
    struct Base {
        virtual ~Base() = default;
        virtual R call(Args...) = 0;
        virtual unique_ptr<Base> clone() = 0;
    };
    
    template<typename F>
    struct Derived : Base {
        F func;
        Derived(F f) : func(f) {}
        R call(Args... args) override {
            return func(args...);
        }
        unique_ptr<Base> clone() override {
            return make_unique<Derived>(*this);
        }
    };
    
    unique_ptr<Base> impl;
    
public:
    template<typename F>
    Function(F f) : impl(make_unique<Derived<F>>(f)) {}
    
    R operator()(Args... args) {
        return impl->call(args...);
    }
    
    Function(const Function& other) : impl(other.impl->clone()) {}
};
\end{lstlisting}

\hypertarget{concept-based-type-erasure-c20}{%
\section{7.3 Concept-Based Type Erasure
(C++20)}\label{concept-based-type-erasure-c20}}

Instead of inheritance, we can erase types that satisfy a concept.

\begin{lstlisting}[language={C++}]
#include <concepts>
#include <memory>

template<typename T>
concept Drawable = requires(T x) { x.draw(); };

class AnyDrawable {
    struct Concept {
        virtual ~Concept() = default;
        virtual void draw() = 0;
        virtual std::unique_ptr<Concept> clone() = 0;
    };

    template<Drawable T>
    struct Model : Concept {
        T data;
        Model(T x) : data(std::move(x)) {}
        void draw() override { data.draw(); }
        std::unique_ptr<Concept> clone() override { return std::make_unique<Model>(data); }
    };

    std::unique_ptr<Concept> pimpl;

public:
    template<Drawable T>
    AnyDrawable(T x) : pimpl(std::make_unique<Model<T>>(std::move(x))) {}
    
    AnyDrawable(const AnyDrawable& other) : pimpl(other.pimpl->clone()) {}
    
    void draw() { pimpl->draw(); }
};
\end{lstlisting}

\begin{center}\rule{0.5\linewidth}{0.5pt}\end{center}

\hypertarget{crtp-curiously-recurring-template-pattern}{%
\section{CRTP (CURIOUSLY RECURRING TEMPLATE
PATTERN)}\label{crtp-curiously-recurring-template-pattern}}

\hypertarget{static-polymorphism}{%
\section{8.1 Static Polymorphism}\label{static-polymorphism}}

\begin{lstlisting}[language={C++}]
// Base class template
template<typename Derived>
class Shape {
public:
    void draw() {
        static_cast<Derived*>(this)->draw_impl();
    }
    
    double area() {
        return static_cast<Derived*>(this)->area_impl();
    }
};

// Derived classes
class Circle : public Shape<Circle> {
private:
    double radius;
    
public:
    Circle(double r) : radius(r) {}
    
    void draw_impl() {
        cout << "Drawing circle\n";
    }
    
    double area_impl() {
        return 3.14159 * radius * radius;
    }
};

class Square : public Shape<Square> {
private:
    double side;
    
public:
    Square(double s) : side(s) {}
    
    void draw_impl() {
        cout << "Drawing square\n";
    }
    
    double area_impl() {
        return side * side;
    }
};

// Usage - no virtual function overhead!
Circle c(5);
c.draw();
cout << c.area() << "\n";

Square s(4);
s.draw();
cout << s.area() << "\n";
\end{lstlisting}

\hypertarget{crtp-for-comparisons}{%
\section{8.2 CRTP for Comparisons}\label{crtp-for-comparisons}}

\begin{lstlisting}[language={C++}]
template<typename Derived>
class Comparable {
public:
    bool operator<(const Comparable& other) const {
        return static_cast<const Derived*>(this)->compare(
            static_cast<const Derived&>(other)
        ) < 0;
    }
    
    bool operator>(const Comparable& other) const {
        return other < *this;
    }
    
    bool operator<=(const Comparable& other) const {
        return !(other < *this);
    }
    
    bool operator>=(const Comparable& other) const {
        return !(*this < other);
    }
    
    bool operator==(const Comparable& other) const {
        return !(*this < other) && !(other < *this);
    }
    
    bool operator!=(const Comparable& other) const {
        return !(*this == other);
    }
};

class Value : public Comparable<Value> {
private:
    int val;
    
public:
    Value(int v) : val(v) {}
    
    int compare(const Value& other) const {
        return val - other.val;
    }
};

// Usage
Value v1(5), v2(10);
cout << (v1 < v2) << "\n";   // true
cout << (v1 == v2) << "\n";  // false
\end{lstlisting}

\begin{center}\rule{0.5\linewidth}{0.5pt}\end{center}

\hypertarget{perfect-forwarding-move-semantics}{%
\section{PERFECT FORWARDING \& MOVE
SEMANTICS}\label{perfect-forwarding-move-semantics}}

\hypertarget{forwarding-problems}{%
\section{9.1 Forwarding Problems}\label{forwarding-problems}}

\begin{lstlisting}[language={C++}]
// The problem: losing information about lvalue/rvalue
template<typename T>
void bad_forward(T arg) {
    // arg is always lvalue
    sink(arg);  // Loses rvalue status
}

// Solution: universal references + std::forward
template<typename T>
void good_forward(T&& arg) {
    // Preserves lvalue/rvalue nature
    sink(std::forward<T>(arg));
}

// Double forwarding
template<typename T>
void wrapper(T&& arg) {
    process(std::forward<T>(arg));
}

template<typename T>
void double_wrapper(T&& arg) {
    wrapper(std::forward<T>(arg));
}
\end{lstlisting}

\hypertarget{move-semantics-implementation}{%
\section{9.2 Move Semantics
Implementation}\label{move-semantics-implementation}}

\begin{lstlisting}[language={C++}]
class String {
private:
    char* data;
    size_t size;
    
public:
    // Move constructor
    String(String&& other) noexcept 
        : data(other.data), size(other.size) {
        other.data = nullptr;
        other.size = 0;
    }
    
    // Move assignment
    String& operator=(String&& other) noexcept {
        if (this != &other) {
            delete[] data;
            data = other.data;
            size = other.size;
            other.data = nullptr;
            other.size = 0;
        }
        return *this;
    }
    
    // Copy constructor (for lvalues)
    String(const String& other)
        : size(other.size), data(new char[size + 1]) {
        strcpy(data, other.data);
    }
    
    // Copy assignment (for lvalues)
    String& operator=(const String& other) {
        if (this != &other) {
            delete[] data;
            size = other.size;
            data = new char[size + 1];
            strcpy(data, other.data);
        }
        return *this;
    }
};
\end{lstlisting}

\begin{center}\rule{0.5\linewidth}{0.5pt}\end{center}

\hypertarget{compile-time-programming}{%
\section{COMPILE-TIME PROGRAMMING}\label{compile-time-programming}}

\hypertarget{tuple-operations-at-compile-time}{%
\section{10.1 Tuple Operations at
Compile-Time}\label{tuple-operations-at-compile-time}}

\begin{lstlisting}[language={C++}]
#include <tuple>
#include <iostream>

// Compile-time tuple iteration
namespace detail {
    template<typename F, typename Tuple, size_t... Is>
    void for_each_impl(F&& f, Tuple&& t, index_sequence<Is...>) {
        (..., f(get<Is>(forward<Tuple>(t))));
    }
}

template<typename F, typename Tuple>
void for_each(F&& f, Tuple&& t) {
    constexpr auto size = tuple_size_v<decay_t<Tuple>>;
    detail::for_each_impl(
        forward<F>(f),
        forward<Tuple>(t),
        make_index_sequence<size>{}
    );
}

// Usage
auto t = make_tuple(1, 2.5, "hello");
for_each([](auto x) { cout << x << " "; }, t);  // 1 2.5 hello
\end{lstlisting}

\hypertarget{compile-time-string-processing}{%
\section{10.2 Compile-Time String
Processing}\label{compile-time-string-processing}}

\begin{lstlisting}[language={C++}]
#include <array>
#include <string_view>

// Compile-time string hashing
constexpr size_t hash_compile_time(string_view s) {
    size_t h = 0;
    for (char c : s) {
        h = h * 31 + c;
    }
    return h;
}

// Compile-time string search
constexpr bool contains_compile_time(string_view s, string_view sub) {
    return s.find(sub) != string_view::npos;
}

// Usage in compile-time context
static_assert(contains_compile_time("hello world", "world"));
constexpr auto hash = hash_compile_time("key");
\end{lstlisting}

\begin{center}\rule{0.5\linewidth}{0.5pt}\end{center}

\hypertarget{meta-object-protocol}{%
\section{META-OBJECT PROTOCOL}\label{meta-object-protocol}}

\hypertarget{reflection-like-patterns}{%
\section{11.1 Reflection-Like Patterns}\label{reflection-like-patterns}}

\begin{lstlisting}[language={C++}]
// Manual reflection without std::meta
struct Member {
    string_view name;
    string_view type;
    size_t offset;
};

template<typename T>
constexpr auto get_members() {
    // Specialize for each type
    return array<Member, 0>{};
}

template<>
constexpr auto get_members<Person>() {
    return array<Member, 3>{{
        {"name", "string", offsetof(Person, name)},
        {"age", "int", offsetof(Person, age)},
        {"email", "string", offsetof(Person, email)}
    }};
}

// Serialize any type
template<typename T>
string serialize(const T& obj) {
    string result;
    for (const auto& member : get_members<T>()) {
        result += member.name;
        result += ": ";
        // Serialize value...
    }
    return result;
}
\end{lstlisting}

\hypertarget{magic-enum-reflection-hack}{%
\section{11.2 Magic Enum (Reflection
Hack)}\label{magic-enum-reflection-hack}}

Before C++26 Reflection, we use compiler intrinsics to get enum names.

\begin{lstlisting}[language={C++}]
template <auto V>
constexpr std::string_view get_enum_name() {
    // Compiler-specific macros
    #ifdef __clang__
        return __PRETTY_FUNCTION__;
    #elif defined(__GNUC__)
        return __PRETTY_FUNCTION__;
    #elif defined(_MSC_VER)
        return __FUNCSIG__;
    #endif
}

enum class Color { Red, Green, Blue };

int main() {
    // Output contains "Color::Red" buried in the string
    std::cout << get_enum_name<Color::Red>() << "\n";
}
\end{lstlisting}

\emph{Note: Libraries like \passthrough{\lstinline!magic\_enum!}
automate parsing this string.}

\begin{center}\rule{0.5\linewidth}{0.5pt}\end{center}

\hypertarget{advanced-container-techniques}{%
\section{ADVANCED CONTAINER
TECHNIQUES}\label{advanced-container-techniques}}

\hypertarget{cow-copy-on-write-string}{%
\section{12.1 COW (Copy-On-Write)
String}\label{cow-copy-on-write-string}}

\begin{lstlisting}[language={C++}]
class CowString {
private:
    struct Buffer : public enable_shared_from_this<Buffer> {
        string data;
        
        Buffer(const string& s) : data(s) {}
    };
    
    shared_ptr<Buffer> buffer;
    
    void ensure_unique() {
        if (!buffer.unique()) {
            buffer = make_shared<Buffer>(*buffer);
        }
    }
    
public:
    CowString(const string& s = "")
        : buffer(make_shared<Buffer>(s)) {}
    
    const char* data() const {
        return buffer->data.c_str();
    }
    
    char& operator[](size_t i) {
        ensure_unique();
        return buffer->data[i];
    }
};
\end{lstlisting}

\hypertarget{intrusive-containers}{%
\section{12.2 Intrusive Containers}\label{intrusive-containers}}

\begin{lstlisting}[language={C++}]
#include <list>

// Node that contains its own link
class IntrusiveNode {
public:
    IntrusiveNode* next = nullptr;
    IntrusiveNode* prev = nullptr;
};

template<typename T>
class IntrusiveList {
private:
    IntrusiveNode* head;
    IntrusiveNode* tail;
    
public:
    void push_back(T* node) {
        if (!head) {
            head = tail = node;
        } else {
            tail->next = node;
            node->prev = tail;
            tail = node;
        }
    }
};
\end{lstlisting}

\begin{center}\rule{0.5\linewidth}{0.5pt}\end{center}

\hypertarget{abi-binary-compatibility}{%
\section{ABI \& BINARY COMPATIBILITY}\label{abi-binary-compatibility}}

\hypertarget{stable-abi-design}{%
\section{13.1 Stable ABI Design}\label{stable-abi-design}}

\begin{lstlisting}[language={C++}]
// Version-stable interface
class StableObject {
private:
    // Use void* for opaque pointer
    void* impl;
    
public:
    StableObject();
    ~StableObject();
    
    // Stable functions
    void do_something(int x);
    int get_value() const;
    
    // Virtual table for future extensions
    virtual ~StableObject() = default;
};

// Implementation in separate library
class StableObjectImpl {
public:
    int value = 0;
    void do_something(int x);
    int get_value() const { return value; }
};

StableObject::StableObject() 
    : impl(new StableObjectImpl()) {}

StableObject::~StableObject() {
    delete static_cast<StableObjectImpl*>(impl);
}

void StableObject::do_something(int x) {
    static_cast<StableObjectImpl*>(impl)->do_something(x);
}

int StableObject::get_value() const {
    return static_cast<StableObjectImpl*>(impl)->get_value();
}
\end{lstlisting}

\begin{center}\rule{0.5\linewidth}{0.5pt}\end{center}

\hypertarget{performance-profiling-optimization}{%
\section{PERFORMANCE PROFILING \&
OPTIMIZATION}\label{performance-profiling-optimization}}

\hypertarget{benchmarking-framework}{%
\section{14.1 Benchmarking Framework}\label{benchmarking-framework}}

\begin{lstlisting}[language={C++}]
#include <chrono>
#include <vector>

class Benchmark {
private:
    using Clock = chrono::high_resolution_clock;
    vector<chrono::nanoseconds> times;
    
public:
    template<typename F>
    void run(F func, int iterations = 1000) {
        for (int i = 0; i < iterations; i++) {
            auto start = Clock::now();
            func();
            auto end = Clock::now();
            times.push_back(chrono::duration_cast<chrono::nanoseconds>(end - start));
        }
    }
    
    void report() {
        if (times.empty()) return;
        
        sort(times.begin(), times.end());
        
        auto sum = 0LL;
        for (auto t : times) sum += t.count();
        
        cout << "Min: " << times.front().count() << " ns\n";
        cout << "Max: " << times.back().count() << " ns\n";
        cout << "Avg: " << (sum / times.size()) << " ns\n";
        cout << "Median: " << times[times.size() / 2].count() << " ns\n";
    }
};

// Usage
Benchmark bm;
bm.run([]() { /* test code */ }, 10000);
bm.report();
\end{lstlisting}

\hypertarget{cache-friendly-algorithms}{%
\section{14.2 Cache-Friendly
Algorithms}\label{cache-friendly-algorithms}}

\begin{lstlisting}[language={C++}]
// Cache-friendly data access pattern
template<typename T>
void cache_friendly_process(vector<T>& data) {
    // Process in cache-line aligned chunks
    const size_t CACHE_LINE = 64;  // 64 bytes typical
    const size_t CHUNK_SIZE = CACHE_LINE / sizeof(T);
    
    for (size_t i = 0; i < data.size(); i += CHUNK_SIZE) {
        // Process one cache line worth of data
        for (size_t j = i; j < min(i + CHUNK_SIZE, data.size()); j++) {
            process_element(data[j]);
        }
    }
}

// SIMD-friendly loop
void simd_process(int* data, size_t size) {
    // Align data for SIMD operations
    alignas(32) int buffer[32];
    
    for (size_t i = 0; i < size; i += 32) {
        // Process 32-byte chunks (8 ints)
        for (int j = 0; j < 8; j++) {
            data[i + j] = process(data[i + j]);
        }
    }
}
\end{lstlisting}

\begin{center}\rule{0.5\linewidth}{0.5pt}\end{center}

\hypertarget{domain-specific-language-design}{%
\section{DOMAIN-SPECIFIC LANGUAGE
DESIGN}\label{domain-specific-language-design}}

\hypertarget{expression-dsl}{%
\section{15.1 Expression DSL}\label{expression-dsl}}

\begin{lstlisting}[language={C++}]
// Domain-specific language for math expressions
class Expr {
public:
    virtual ~Expr() = default;
    virtual double evaluate() = 0;
};

class Constant : public Expr {
private:
    double value;
public:
    Constant(double v) : value(v) {}
    double evaluate() override { return value; }
};

class BinaryOp : public Expr {
private:
    shared_ptr<Expr> left, right;
    function<double(double, double)> op;
    
public:
    BinaryOp(shared_ptr<Expr> l, shared_ptr<Expr> r,
             function<double(double, double)> op)
        : left(l), right(r), op(op) {}
    
    double evaluate() override {
        return op(left->evaluate(), right->evaluate());
    }
};

// Operator overloading for DSL
auto operator+(shared_ptr<Expr> l, shared_ptr<Expr> r) {
    return make_shared<BinaryOp>(l, r, [](double a, double b) { return a + b; });
}

// Usage
auto expr = make_shared<Constant>(3) + make_shared<Constant>(4);
cout << expr->evaluate() << "\n";  // 7
\end{lstlisting}

\begin{center}\rule{0.5\linewidth}{0.5pt}\end{center}

\hypertarget{modern-design-patterns}{%
\section{MODERN DESIGN PATTERNS}\label{modern-design-patterns}}

Traditional GoF patterns often use inheritance. Modern C++ favors
composition, templates, and lambdas.

\hypertarget{strategy-pattern-functional-approach}{%
\section{16.1 Strategy Pattern (Functional
Approach)}\label{strategy-pattern-functional-approach}}

Instead of a class hierarchy, use
\passthrough{\lstinline!std::function!} or templates.

\begin{lstlisting}[language={C++}]
#include <functional>
#include <iostream>
#include <vector>

// Traditional: abstract base class Strategy
// Modern: std::function
using SortStrategy = std::function<void(std::vector<int>&)>;

class Sorter {
    SortStrategy strategy;
public:
    Sorter(SortStrategy s) : strategy(s) {}
    
    void sort(std::vector<int>& data) {
        if (strategy) strategy(data);
    }
};

int main() {
    std::vector<int> data = {5, 2, 9, 1};
    
    // Strategy 1: Lambda
    Sorter s1([](auto& v) { std::sort(v.begin(), v.end()); });
    
    // Strategy 2: Different logic
    Sorter s2([](auto& v) { std::sort(v.rbegin(), v.rend()); });
    
    s1.sort(data);
    return 0;
}
\end{lstlisting}

\hypertarget{visitor-pattern-stdvariant}{%
\section{16.2 Visitor Pattern
(std::variant)}\label{visitor-pattern-stdvariant}}

Replace virtual functions with \passthrough{\lstinline!std::variant!}
and \passthrough{\lstinline!std::visit!} for closed sets of types.

\begin{lstlisting}[language={C++}]
#include <variant>
#include <iostream>
#include <vector>

struct Circle { double radius; };
struct Square { double side; };
using Shape = std::variant<Circle, Square>;

// Visitor
struct AreaVisitor {
    double operator()(const Circle& c) { return 3.14159 * c.radius * c.radius; }
    double operator()(const Square& s) { return s.side * s.side; }
};

int main() {
    std::vector<Shape> shapes = { Circle{2.0}, Square{3.0} };
    
    for (const auto& s : shapes) {
        // Apply visitor
        double area = std::visit(AreaVisitor{}, s);
        std::cout << "Area: " << area << "\n";
    }
    
    // With lambda (overloaded pattern)
    // See "Helper for std::visit" in many codebases
    return 0;
}
\end{lstlisting}

\begin{center}\rule{0.5\linewidth}{0.5pt}\end{center}

\hypertarget{hardware-sympathy}{%
\section{HARDWARE SYMPATHY}\label{hardware-sympathy}}

\hypertarget{cache-locality-false-sharing}{%
\section{17.1 Cache Locality \& False
Sharing}\label{cache-locality-false-sharing}}

CPUs load data in cache lines (typically 64 bytes).

\hypertarget{false-sharing}{%
\subsection{False Sharing}\label{false-sharing}}

When two threads modify independent variables that sit on the
\emph{same} cache line, they invalidate each other's cache, destroying
performance.

\begin{lstlisting}[language={C++}]
#include <new>
#include <atomic>

struct BadCounter {
    std::atomic<int> a; // Thread 1 modifies
    std::atomic<int> b; // Thread 2 modifies
    // Likely on same cache line -> ping-pong effect
};

struct GoodCounter {
    alignas(64) std::atomic<int> a; // Forced to own cache line
    alignas(64) std::atomic<int> b;
};
\end{lstlisting}

\hypertarget{branch-prediction}{%
\section{17.2 Branch Prediction}\label{branch-prediction}}

CPUs try to guess which way an \passthrough{\lstinline!if!} will go.
Modern C++20 provides attributes to help.

\begin{lstlisting}[language={C++}]
void process(int* ptr) {
    if (!ptr) [[unlikely]] {
        // Compiler optimizes this block to be "cold"
        // CPU assumes this won't happen
        throw std::runtime_error("Null pointer");
    }
    
    // This "hot" path is optimized for fall-through
    if (ptr) [[likely]] {
        *ptr = 42;
    }
}
\end{lstlisting}

\hypertarget{simd-single-instruction-multiple-data}{%
\section{17.3 SIMD (Single Instruction, Multiple
Data)}\label{simd-single-instruction-multiple-data}}

Using intrinsics (or libraries like \passthrough{\lstinline!std::simd!}
in future) to process data in parallel lanes.

\begin{lstlisting}[language={C++}]
// Example: Manual unrolling for auto-vectorization
void add_arrays(float* a, float* b, float* c, int n) {
    // Tell compiler pointers don't alias (C99 restrict, or implementation specific)
    // #pragma omp simd 
    for (int i = 0; i < n; ++i) {
        c[i] = a[i] + b[i];
    }
}
\end{lstlisting}

\begin{center}\rule{0.5\linewidth}{0.5pt}\end{center}
